
\subsubsection{Interpretation}

The results suggest that cross-dimensional effects exist but exhibit heterogeneous patterns. LoRA interventions targeting attention layers caused representation changes throughout the network (100\% layer coverage), with larger magnitude changes in attention mechanisms (2× compared to residual streams). Despite universal representation changes, dimension pairs showed varying correlation patterns.

The near-zero truthfulness-fairness correlation (r = 0.034) suggests these dimensions may operate through relatively independent mechanisms. The strong negative truthfulness-robustness correlation observed in GPT-2 (r = -0.997, though p = 0.051) suggests potential competition, but this pattern did not replicate consistently across architectures. The fairness-robustness relationship showed negative correlation direction in two of three models, with statistical significance in one (OPT: r = -0.886, p = 0.046).

One interpretation is that some trustworthiness dimensions share computational resources (fairness and robustness, which both involve contextual reasoning), while others rely on more independent mechanisms (truthfulness may primarily involve factual knowledge retrieval separate from bias suppression). However, this interpretation should be considered tentative given the limited statistical power and partial replication across architectures.

\subsubsection{Limitations}

Several limitations constrain interpretation of these findings:

\textbf{Statistical power}: Experiments used 3-5 replicates per condition, providing limited power to detect effects. The truthfulness-robustness correlation (r = -0.997, p = 0.051) exemplifies this: large effect size but marginal statistical significance. Larger sample sizes (n ≥ 10) would enable more precise estimation.

\textbf{Architecture coverage}: Only transformer variants were tested (GPT-2, OPT, Pythia). Claims of generalization apply only within the transformer family. Non-transformer architectures (state-space models, recurrent networks) may exhibit different patterns.

\textbf{Intervention method}: All experiments used LoRA targeting attention layers. Full fine-tuning, prompt-based methods, or interventions targeting different components might produce different cross-dimensional effects. The observed patterns may be LoRA-specific rather than general properties of trustworthiness interventions.

\textbf{Dimension coverage}: Three dimensions were evaluated (truthfulness, fairness, robustness), omitting privacy, safety, and machine ethics. The observed correlation patterns may not extend to unmeasured dimensions.

\textbf{Benchmark evaluation}: Performance was measured on established benchmarks (TruthfulQA, BBQ, ANLI), which serve as proxies for trustworthiness but may not fully capture real-world behavior. Correlation patterns observed on benchmarks may differ from relationships in deployment contexts.

\textbf{Unexpected performance changes}: In H-M3, truthfulness decreased post-intervention (26.0\% vs. 29.4\% baseline) despite targeted training on truthfulness data. This may reflect overfitting to the training subset, distribution shift, or evaluation artifacts. This result complicates interpretation of cross-dimensional effects.

\subsubsection{Implications}

If the observed patterns replicate with adequate statistical power, they would suggest that practitioners cannot assume independence when applying targeted trustworthiness interventions. Some dimension pairs may require explicit multi-objective optimization to avoid unintended trade-offs. However, current evidence is insufficient to make strong practical recommendations given the limitations noted above.

The perturbation-based methodology demonstrates one approach to characterizing cross-dimensional relationships by generating controlled variance through systematic parameter variation. This approach could be applied to other intervention types, architectures, or dimension combinations.
